\documentclass[11pt]{article}
\usepackage{amsmath, amssymb, amsthm}
\usepackage{graphicx}
\usepackage{draftwatermark}
\SetWatermarkText{SUPERSEDED}
\SetWatermarkScale{1.5}
\SetWatermarkAngle{45}
\SetWatermarkColor[gray]{0.80}

\newtheorem{definition}{Definition}
\newtheorem{lemma}{Lemma}[section]
\newtheorem{proposition}{Proposition}
\newtheorem{postulate}{Postulate}
\newtheorem{conjecture}{Conjecture}
\newtheorem{corollary}{Corollary}
\newtheorem{remark}{Remark}

\title{A Regular-Language and Tree Representation of Odd Collatz Dynamics}
\author{Jon Seymour\\
  \small\texttt{jon@wildducktheories.com}\\
  \small Sydney, Australia}
\date{June 2026}

\begin{document}
\maketitle

\begin{abstract}
We study the odd Collatz map $T(n) = (3n+1)/2^{v_2(3n+1)}$ by classifying
each step according to the 2-adic valuation $v = v_2(3n+1)$: an OE step
($v=1$, $n\equiv 3$ or $7\pmod{8}$), an OEE step ($v=2$, $n\equiv
1\pmod{8}$), or an $\mathrm{OEEE}^+$ step ($v\geq 3$, $n\equiv 5\pmod{8}$).
We show that every Collatz orbit is partitioned into \emph{Steiner circuits},
each matching the regular expression $(7^*3)?(1\mid 5)$ over the mod-8
residues $\{1,3,5,7\}$, with full orbits matching $((7^*3)?(1\mid 5))^*$.
We further introduce the \emph{$5\pmod{8}$ overlay tree} for a forward
Collatz path from an initial node $a$ to a final node $b$: its nodes are
$a$, $b$, all $5\pmod{8}$ nodes on the path, and for each such node the
$0\pmod{3}$ node $m$ such that the forward OEEE-free path from $m$ reaches $n$.
We illustrate the construction on the path $27\to 1$, which is a
concatenation of several OEEE-free paths joined at $5\pmod{8}$ nodes.
For any two odd integers $a$, $b$ connected by a Collatz path with $o\geq 1$
steps, we prove $2^e < 3^o \implies a < b$; for OEEE-free pairs we prove
the biconditional $2^e \geq 3^o \iff a \geq b$ conditional on
Conjecture~3 (OE negative parity), via an inductive argument on the ordering
target function $\Delta = (3^o - 2^e)a + K$; the OEE cases are proved
unconditionally and the conjecture precisely identifies the remaining gap.  All 289 failures of the ordering implication found across
an adversarial sample of 240{,}000 pairs occurred on paths containing at least
one interior node $\equiv 5\pmod{8}$, with none on the 57{,}381 OEEE-free pairs
examined, confirming that the $5\pmod{8}$ condition is the essential mechanism
by which the ordering can fail.  We prove that the only all-OEE cycle is $1\to 1$ and conjecture
this extends to all OEEE-free cycles; hence any other Collatz cycle must be an
OEEE-path.  We further conjecture that no Collatz cycle contains a $5\pmod{8}$
node; conditional on this and the unique-cycle conjecture, there are no Collatz
cycles other than $1\to 1$.
\end{abstract}

\section{The Collatz Map and Step Classification}

We work with odd positive integers under the Collatz map.  For an odd integer
$n$, define
\[
  T(n) \;=\; \frac{3n+1}{2^{v_2(3n+1)}},
\]
where $v_2(m)$ denotes the 2-adic valuation of $m$ (the largest power of $2$
dividing $m$).  By construction $T(n)$ is always odd.

\begin{definition}[OE, OEE, OEEE$^+$, and OEE$^+$ steps]
Let $n$ be an odd positive integer and set $v = v_2(3n+1)$.
\begin{itemize}
  \item An \textbf{OE step} occurs when $v = 1$, i.e.\ $3n+1 \equiv 2
        \pmod{4}$, equivalently $n \equiv 3 \pmod{4}$, i.e.\ $n \equiv 3$
        or $7 \pmod{8}$.  Interior OE steps (those not immediately preceding
        an OEE$^+$ step) have $n \equiv 7\pmod{8}$; the unique final OE
        step of each OE-run, which immediately precedes the OEE$^+$
        termination, has $n \equiv 3\pmod{8}$.
  \item An \textbf{OEE step} occurs when $v = 2$, i.e.\ $3n+1 \equiv 4
        \pmod{8}$, equivalently $n \equiv 1 \pmod{8}$.
  \item An $\mathbf{OEEE}^+$ \textbf{step} occurs when $v \geq 3$, i.e.\
        $3n+1 \equiv 0 \pmod{8}$, equivalently $n \equiv 5 \pmod{8}$ (when
        $v=3$) or higher powers of $2$ divide $3n+1$.
  \item An $\mathbf{OEE}^+$ \textbf{step} is either an OEE or an
        $\mathrm{OEEE}^+$ step, i.e.\ $v \geq 2$; it corresponds to
        residue $1$ or $5 \pmod 8$.
\end{itemize}
\end{definition}

\begin{definition}[Steiner circuit]
A \textbf{Steiner circuit} is a maximal run of Collatz steps of the form
\[
  (OE)^*\,(OEE^+),
\]
that is, a (possibly empty) sequence of OE steps followed by a single
terminating $OEE^+$ step (which is either OEE or $\mathrm{OEEE}^+$).
\end{definition}

\begin{definition}[OEEE-free pair, OEEE-free path, OEEE-path]
\label{def:branches}
Let $a$ and $b$ be odd positive integers with $b$ reachable from $a$
under repeated application of $T$.  Write $T^{(k)}(a)$ for the $k$-th
iterate of $T$ starting at $a$, with $T^{(0)}(a) = a$, and let $o \geq 1$
be the number of steps from $a$ to $b$, so that $T^{(o)}(a) = b$.
We call $(a, b)$ an \textbf{OEEE-free pair} if
\[
  \bigl\{\, T^{(k)}(a) \;\mid\; 0 \leq k < o \,\bigr\}
  \;\cap\; \{\,n \in \mathbb{Z}_{>0} \mid n \equiv 5\pmod{8}\,\}
  \;=\; \emptyset,
\]
i.e.\ no iterate strictly before $b$ is $\equiv 5\pmod{8}$; the endpoint
$b$ itself may or may not be $\equiv 5\pmod{8}$.  The forward path from
$a$ to $b$ is then called an \textbf{OEEE-free path}; otherwise it is an
\textbf{OEEE-path}.
\end{definition}

\medskip
We use \textbf{path} exclusively for forward traversals of the Collatz map
(from a node toward $1$) and \textbf{walk} for traversals in the reverse
direction (from a node toward its predecessors, i.e.\ from root toward
leaves in the Collatz tree).

\begin{definition}[OEEE-free walk]
\label{def:walk}
An \textbf{OEEE-free walk} is a sequence of nodes $n_0, n_1, n_2, \ldots$
governed by the following state machine on $n \bmod 3$:
\[
\begin{array}{cll}
n \bmod 3 & \text{action} & \text{next node} \\
\hline
0 & \textbf{halt} & - \\
1 & \text{multiply-shift} & (4n-1)/3 \\
2 & \text{multiply-shift} & (2n-1)/3 \\
\end{array}
\]
Each multiply-shift step corresponds to $k \leq 2$ halvings in the forward
Collatz map ($k=1$ for $n\equiv 2\pmod{3}$; $k=2$ for $n\equiv 1\pmod{3}$),
so every step has $v \leq 2$ and the walk is OEEE-free by construction.
The walk halts at the first node $\equiv 0\pmod{3}$; it is postulated
to always halt (Postulate~\ref{post:halt}).  If a walk starts at $b$ and
halts at $a$, then there exists a forward Collatz path from $a$ to $b$.
\end{definition}

\subsection*{A note on exponential versus logarithmic form}

Throughout this paper we state ordering conditions in the exponential
form $2^e \geq 3^o$ rather than the equivalent logarithmic form
$e \geq o\log_2 3$.  The two are equivalent:

\begin{lemma}[Exponential--logarithmic equivalence]
\label{lem:explog}
For non-negative integers $e$ and $o$,
\[
  2^e \;\geq\; 3^o \;\iff\; e \;\geq\; o\log_2 3.
\]
\end{lemma}

\begin{proof}
Since $x \mapsto 2^x$ is strictly increasing and $2^{o\log_2 3} = 3^o$,
we have $2^e \geq 3^o \iff 2^e \geq 2^{o\log_2 3} \iff e \geq o\log_2 3$.
\end{proof}

\noindent The logarithmic form $e \geq o\log_2 3$ has intuitive appeal as
a comparison of step-count ratios against $\log_2 3 \approx 1.585$, and
we retain it in Section~\ref{sec:approx} for motivational purposes.
However, $\log_2 3$ is transcendental (a consequence of the
Gel'fond--Schneider theorem), so any inequality involving it carries an
implicit appeal to real analysis.  The exponential form $2^e \geq 3^o$
involves only integers and is decidable by direct comparison.  This
preference for exact integer arithmetic over analytic approximation is
in the spirit of Erd\H{o}s's proof of Bertrand's postulate~\cite{erdos1932},
which replaced analytic estimates with explicit binomial coefficient
bounds.  All theorems and conjectures in this paper are therefore stated
in the exponential form, even though the two forms are logically
interchangeable by Lemma~\ref{lem:explog}.

\section{Mod 8 Encoding and the Regular Language}

\subsection{Why each step type carries a fixed residue mod 8}

We track residues of successive iterates modulo 8.

\begin{proposition}[Mod-8 residues of OE/OEE/$\mathrm{OEEE}^+$ steps]
\label{prop:mod8}
The step type of each odd iterate $n$ is determined by $n \bmod 8$:
\begin{itemize}
  \item $n \equiv 7 \pmod 8$: $3n+1 \equiv 6 \pmod{16}$, so
        $v_2(3n+1) = 1$ (\textbf{OE step}).
  \item $n \equiv 3 \pmod 8$: $3n+1 \equiv 10 \pmod{16}$, so
        $v_2(3n+1) = 1$ (\textbf{OE step}).  This is the \emph{final} OE
        step of an OE-run, since $T(n) = (3n+1)/2 \equiv 5\pmod{8}$ is
        itself an OEEE$^+$ node.
  \item $n \equiv 1 \pmod 8$: $3n+1 \equiv 4 \pmod{16}$, so
        $v_2(3n+1) = 2$ (\textbf{OEE step}).
  \item $n \equiv 5 \pmod 8$: $3n+1 \equiv 0 \pmod{8}$, so
        $v_2(3n+1) \geq 3$ ($\mathbf{OEEE}^+$ \textbf{step}); the exact
        value of $v_2(3n+1)$ and of $T(n) \bmod 8$ depend on higher bits
        of $n$.
\end{itemize}
Note that $T(n) \bmod 8$ is \emph{not} determined by $n \bmod 8$ alone for
OE and OEE steps; the encoding uses only the \emph{input} residue $n\bmod 8$.
\end{proposition}

The key point is that the \emph{input} residue $n \bmod 8$ determines the
symbol we assign in the encoding, independently of what happens after the
even divisions.

For the purpose of symbolic encoding we record the \emph{residue class of the
input node} $n$ within each step:

\begin{center}
\begin{tabular}{llc}
\hline
Residue $n \bmod 8$ & Step type & Symbol \\
\hline
$7$ & OE (interior of run) & $7$ \\
$3$ & OE (final step before OEE$^+$ termination) & $3$ \\
$1$ & OEE (exactly two even divisions) & $1$ \\
$5$ & OEEE$^+$ (three or more even divisions) & $5$ \\
\hline
\end{tabular}
\end{center}

\subsection{The regular expression}

The correspondence above implies a definite structure:

\begin{itemize}
  \item A maximal \textbf{OE-run} starts with zero or more residue-$7$ nodes
        and ends at a residue-$3$ node; it is encoded as $7^*3$.
  \item The symbol \textbf{3} marks the \emph{final OE step before the
        OEE$^+$ termination} of the current OE-run, not the end of the
        Steiner circuit itself.
  \item Termination by \textbf{OEE} (exactly two even divisions) occurs at a
        node congruent to $1 \pmod 8$, encoded as $1$.
  \item Termination by $\mathbf{OEEE}^+$ (three or more even divisions) occurs
        at a node congruent to $5 \pmod 8$, encoded as $5$.
  \item Together, OEE and $\mathrm{OEEE}^+$ constitute an $\mathbf{OEE}^+$
        termination (two or more even divisions), encoded as $1 \mid 5$.
\end{itemize}

Thus a single Steiner circuit---zero or more OE-runs followed by a terminating
$\mathrm{OEE}^+$ step (OEE or $\mathrm{OEEE}^+$)---is encoded by the regular expression:
\[
  \boxed{(7^*3)?(1 \mid 5)}.
\]

An entire odd Collatz orbit, decomposed into successive Steiner circuits, is
encoded by:
\[
  \boxed{\bigl((7^*3)?(1 \mid 5)\bigr)^*}.
\]

This is a regular language over the alphabet $\{1, 3, 5, 7\}$, giving a
finite-state symbolic grammar for the mod-8 projection of odd Collatz
dynamics.

\section{Example: the odd path $27 \to 1$}

Table~\ref{tab:path} shows the full odd Collatz path from $27$ to $1$
under the map $T$.  The residue $n\bmod 8$ serves directly as the
encoding symbol (Proposition~\ref{prop:mod8}), so no separate symbol
column is needed.

\begin{table}[h]
\centering
\small
\begin{minipage}[t]{0.48\textwidth}
\centering
\begin{tabular}{rrrrc}
\hline
Step & $n$ & $n{\bmod}8$ & $v$ & Type \\
\hline
0  & 27   & 3 & 1 & OE       \\
1  & 41   & 1 & 2 & OEE      \\
2  & 31   & 7 & 1 & OE       \\
3  & 47   & 7 & 1 & OE       \\
4  & 71   & 7 & 1 & OE       \\
5  & 107  & 3 & 1 & OE       \\
6  & 161  & 1 & 2 & OEE      \\
7  & 121  & 1 & 2 & OEE      \\
8  & 91   & 3 & 1 & OE       \\
9  & 137  & 1 & 2 & OEE      \\
10 & 103  & 7 & 1 & OE       \\
11 & 155  & 3 & 1 & OE       \\
12 & 233  & 1 & 2 & OEE      \\
13 & 175  & 7 & 1 & OE       \\
14 & 263  & 7 & 1 & OE       \\
15 & 395  & 3 & 1 & OE       \\
16 & 593  & 1 & 2 & OEE      \\
17 & 445  & 5 & 3 & OEEE$^+$ \\
18 & 167  & 7 & 1 & OE       \\
19 & 251  & 3 & 1 & OE       \\
20 & 377  & 1 & 2 & OEE      \\
\hline
\end{tabular}
\end{minipage}
\hfill
\begin{minipage}[t]{0.48\textwidth}
\centering
\begin{tabular}{rrrrc}
\hline
Step & $n$ & $n{\bmod}8$ & $v$ & Type \\
\hline
21 & 283  & 3 & 1 & OE       \\
22 & 425  & 1 & 2 & OEE      \\
23 & 319  & 7 & 1 & OE       \\
24 & 479  & 7 & 1 & OE       \\
25 & 719  & 7 & 1 & OE       \\
26 & 1079 & 7 & 1 & OE       \\
27 & 1619 & 3 & 1 & OE       \\
28 & 2429 & 5 & 3 & OEEE$^+$ \\
29 & 911  & 7 & 1 & OE       \\
30 & 1367 & 7 & 1 & OE       \\
31 & 2051 & 3 & 1 & OE       \\
32 & 3077 & 5 & 4 & OEEE$^+$ \\
33 & 577  & 1 & 2 & OEE      \\
34 & 433  & 1 & 2 & OEE      \\
35 & 325  & 5 & 4 & OEEE$^+$ \\
36 & 61   & 5 & 3 & OEEE$^+$ \\
37 & 23   & 7 & 1 & OE       \\
38 & 35   & 3 & 1 & OE       \\
39 & 53   & 5 & 5 & OEEE$^+$ \\
40 & 5    & 5 & 4 & OEEE$^+$ \\
41 & 1    & 1 & — & terminal \\
\hline
\end{tabular}
\end{minipage}
\caption{Odd Collatz path $27\to 1$: 41 steps, 42 nodes.  The symbol
  for each step equals $n\bmod 8$.}
\label{tab:path}
\end{table}

The 41-symbol string over $\{1,3,5,7\}$, with Steiner circuits marked by
vertical bars, is:
\begin{align*}
&\underbrace{3\,1}_{1}\mid
 \underbrace{7\,7\,7\,3\,1}_{2}\mid
 \underbrace{1}_{3}\mid
 \underbrace{3\,1}_{4}\mid
 \underbrace{7\,3\,1}_{5}\mid
 \underbrace{7\,7\,3\,1}_{6}\mid
 \underbrace{5}_{7}\mid
 \underbrace{7\,3\,1}_{8}\mid
 \underbrace{3\,1}_{9}\\
&\mid
 \underbrace{7\,7\,7\,7\,3\,5}_{10}\mid
 \underbrace{7\,7\,3\,5}_{11}\mid
 \underbrace{1}_{12}\mid
 \underbrace{1}_{13}\mid
 \underbrace{5}_{14}\mid
 \underbrace{5}_{15}\mid
 \underbrace{7\,3\,5}_{16}\mid
 \underbrace{5}_{17}
\end{align*}
Each group is one Steiner circuit matching $(7^*3)?(1\mid 5)$.
The path comprises 17 Steiner circuits.

\section{The $5\pmod{8}$ Overlay Tree}

\subsection{Definition}

Let $a$ be an odd positive integer (the \textbf{initial node}) and $b$ an
odd positive integer (the \textbf{final node}, typically $b=1$).  The
\textbf{$5\pmod{8}$ overlay tree} with initial node $a$ and final node $b$ is
defined if and only if $b = T^k(a)$ for some $k \geq 0$, i.e.\ $b$ lies
on the forward Collatz path from $a$.

\medskip
\noindent\textbf{Nodes.} The tree always contains:
\begin{itemize}
  \item the initial node $a$,
  \item the final node $b$,
  \item every node $\equiv 5\pmod{8}$ on the forward path from $a$ to $b$,
  \item for each such $5\pmod{8}$ node $n$: the $0\pmod{3}$ node $m$
        reached by the OEEE-free walk starting from $n$ (Definition~\ref{def:walk}),
        if it exists (assumed by Postulate~\ref{post:halt}).
\end{itemize}

\noindent\textbf{Edges.} All edges are directed in the forward Collatz direction
(from $a$ toward $b$).  Each node type carries a specific outgoing edge:
\begin{itemize}
  \item The \textbf{initial node} $a$ points to the first $5\pmod{8}$ node
        on the forward path from $a$ to $b$, or directly to $b$ if no
        $5\pmod{8}$ node lies between $a$ and $b$.
  \item Each $\mathbf{5\pmod{8}}$ \textbf{node} points to the next
        $5\pmod{8}$ node on the forward path toward $b$, or directly to
        $b$ if no further $5\pmod{8}$ node exists between it and $b$.
  \item Each $\mathbf{0\pmod{3}}$ \textbf{node} $m$ points to the first
        $5\pmod{8}$ node on the forward Collatz path from $m$.  Such a node
        $m$ is the terminus of the OEEE-free walk starting from the
        corresponding $5\pmod{8}$ node $n$.
\end{itemize}

\newtheorem{theorem}{Theorem}

\begin{postulate}
\label{post:halt}
The OEEE-free walk (Definition~\ref{def:walk}) always halts, i.e.\ for every
node $n \equiv 5\pmod{8}$, the walk starting from $n$ eventually reaches a
node $\equiv 0\pmod{3}$.  Equivalently, for every node $n \equiv 5\pmod{8}$,
there exists a node $m \equiv 0\pmod{3}$ such that the forward Collatz path
from $m$ reaches $n$ as an OEEE-free path.
\end{postulate}

This is supported by computation but is not proved in general.

\subsection{Example: $a = 27$, $b = 1$}

Figure~\ref{fig:reverse-tree} shows the $5\pmod{8}$ overlay tree with initial
node $a = 27$ and final node $b = 1$.  The forward path from $27$ to $1$
passes through seven nodes $\equiv 5\pmod{8}$, namely $445, 2429, 3077,
325, 61, 53, 5$ (in forward order, $a$ to $b$).  Each of these points to
the next $5\pmod{8}$ node in the forward direction, with the last one ($5$)
pointing to $b=1$.  All seven have $0\pmod{3}$ termini, consistent with
Postulate~\ref{post:halt}.  The initial node $a=27$ points to the first
$5\pmod{8}$ node $445$.  For each $5\pmod{8}$ node $n$, the $0\pmod{3}$
node $m$ whose forward OEEE-free path reaches $n$ has path length ranging
from $1$ step ($n=5$, $n=61$) to $17$ steps ($n=445$, where $m=27=a$).

\begin{figure}[h]
  \centering
  \includegraphics[width=0.95\textwidth]{reverse-tree.pdf}
  \caption{$5\pmod{8}$ overlay tree with initial node $a=27$ (pink) and
    final node $b=1$ (yellow).  Blue nodes are $\equiv 5\pmod{8}$; green
    nodes are $0\pmod{3}$ termini.  All arrows point in the forward Collatz
    direction.  Each $5\pmod{8}$ node points to the next $5\pmod{8}$ node
    on the forward path, or to $b$ if none exists; $a=27$ points to the
    first $5\pmod{8}$ node $445$.  Each $0\pmod{3}$ node points to the
    first $5\pmod{8}$ node on its forward path.
    By definition, no OEEE-free path visits a $5\pmod{8}$ node.}
  \label{fig:reverse-tree}
\end{figure}

\section{The Path Identity and the $3^o/2^e$ Growth Factor}

\subsection{Setup}

Let $a$ and $b$ be odd positive integers such that $b$ is reachable from
$a$ on the forward Collatz path and no node strictly between $a$ and $b$
is $\equiv 5\pmod{8}$; $b$ itself may or may not be $\equiv 5\pmod{8}$.
(The case $b = a$ is the null path with $e = o = 0$.)  Define:
\begin{itemize}
  \item $e$ = the total number of even halvings on the forward path from
        $a$ to $b$, not counting $b$ itself, and
  \item $o$ = the number of odd Collatz steps on the forward path from $a$
        to $b$, not counting $b$ itself.
\end{itemize}

\noindent\textbf{Example.} Take $a = 75$.  The forward path is
\[
  75 \;\to\; 226 \;\to\; 113 \;\to\; 340 \;\to\; 170 \;\to\; 85 = b,
\]
giving $o = 2$ odd steps (at $75, 113$) and $e = 3$ even halvings
(at $226, 340, 170$), with $b = 85$.

\subsection{The Path Identity}

Each odd step multiplies by $3$ and adds $1$; each even step divides by $2$.
Tracking the accumulated powers along the forward path from $a$ to $b$ yields
the \textbf{path identity}:
\[
  2^e \cdot b \;=\; 3^o \cdot a \;+\; K,
\]
where $K \geq 0$ is an integer that collects the carry terms from each
$3n+1$ step, weighted by the remaining powers of $2$.  Dividing through by
$2^e a$:
\[
  \frac{b}{a} \;=\; \frac{3^o}{2^e} + \frac{K}{2^e \cdot a}.
\]
Since $K \geq 0$, the growth factor $b/a$ always exceeds the bare ratio
$3^o/2^e$; the correction $K/(2^e a)$ is non-negative.

\subsection{The Growth Factor $b/a \approx 3^o/2^e$}
\label{sec:approx}

For large $a$ and $b$ the correction term $K/(2^e a)$ is small, giving
\[
  \frac{b}{a} \;\approx\; \frac{3^o}{2^e}.
\]
This can seem surprising: why should the growth factor of a Collatz path
mirror the ratio of powers of $3$ and $2$ raised to the counts of odd
and even steps?  The path identity shows it is in fact an exact algebraic
consequence, with the correction vanishing as $a \to \infty$.  Since
$K \geq 0$ the growth factor always satisfies $b/a \geq 3^o/2^e$; the
correction is always non-negative.

Taking logarithms base $2$:
\[
  \log_2(b/a) \;\approx\; o\log_2 3 - e.
\]
Thus $b > a$ when $3^o > 2^e$ (a surplus of odd steps relative to
even halvings, so the path grows) and $b < a$ when $3^o < 2^e$ (a
deficit of odd steps, so the path contracts).  This makes structural
sense: a surplus of odd steps inflates the iterate relative to where
it started.

\subsection{Sharpness on OEEE-free Paths}

Let $a$ and $b$ be odd positive integers connected by an OEEE-free path.
Since $e$ and $o$ count operations (steps), not nodes, $T$ is never applied
to the final node $b$, so $b$ contributes nothing to $e$ or $o$.  Every
step source on an OEEE-free path is not $\equiv 5\pmod{8}$, so every step
has $v \leq 2$, giving $e \leq 2o$.

\begin{proposition}[OEEE-free carry bounds]
\label{prop:Kbounds}
Let $(a, b)$ be an OEEE-free pair with $o \geq 1$ steps.  Then:
\[
  3^o - 2^o \;\leq\; K \;\leq\; 4^o - 3^o.
\]
Both bounds are tight: the lower bound is achieved when every step is an OE
step ($v = 1$); the upper bound when every step is an OEE step ($v = 2$).
\end{proposition}

\begin{proof}
We use the suffix recurrence: $R_o = 0$ and $R_m = 2^{v_m} R_{m+1} + 3^{o-1-m}$
for $0 \leq m < o$, where $v_m \in \{1, 2\}$ since the path is OEEE-free.
Then $K = R_0$.

\emph{Base case} ($o = 1$): The single step has $v_0 \in \{1, 2\}$ and
$K = R_0 = 2^{v_0} \cdot 0 + 3^0 = 1$.  Since $3^1 - 2^1 = 1 = 4^1 - 3^1$,
the bounds hold with equality.

\emph{Inductive step}: Suppose the bounds hold for a path of length $o$,
i.e.\ $3^o - 2^o \leq K' \leq 4^o - 3^o$.  Prepending one step with
valuation $v \in \{1, 2\}$ gives a path of length $o + 1$ with carry
\[
  K \;=\; 2^v K' + 3^o.
\]
For the lower bound, $v \geq 1$, so
\[
  K \;\geq\; 2(3^o - 2^o) + 3^o \;=\; 3 \cdot 3^o - 2^{o+1} \;=\; 3^{o+1} - 2^{o+1}.
\]
For the upper bound, $v \leq 2$, so
\[
  K \;\leq\; 4(4^o - 3^o) + 3^o \;=\; 4^{o+1} - 3 \cdot 3^o \;=\; 4^{o+1} - 3^{o+1}.
\]
Both inequalities are equalities when every step has $v = 1$ (all-OE) or
$v = 2$ (all-OEE) respectively, confirming tightness.
\end{proof}

Hence on any OEEE-free path:
\[
  3^o - 2^o \;\leq\; K \;\leq\; 4^o - 3^o \;<\; 4^o.
\]

\begin{theorem}[Collatz path ordering, positive parity state]
\label{thm:approx}
Let $a$ and $b$ be odd positive integers connected by a Collatz path
with $o \geq 1$, and let $e$, $o$, $K$ be as above.  Then:
\[
  2^e < 3^o \;\implies\; a < b.
\]
\end{theorem}

\begin{proof}
From the growth-factor form $b/a = 3^o/2^e + K/(2^e a)$ with $K \geq 0$:
if $3^o > 2^e$ then $b/a \geq 3^o/2^e > 1$, so $b > a$.

The approximation bound $|b/a - 3^o/2^e| < (4/3)^o/a$ follows from the
OEEE-free bound $K \leq 4^o - 3^o < 4^o$ of Proposition~\ref{prop:Kbounds}.
\end{proof}

\section{Ordering on OEEE-free Paths}
\label{sec:ordering}

\subsection{The unified ordering target function}

For an OEEE-free path $a \to b$ with parameters $(o, e, K)$, the path
identity (Section~\ref{sec:approx}) gives $2^e b = 3^o a + K$.
Subtracting $2^e a$ from both sides isolates the spatial displacement:
\[
  2^e(b - a) \;=\; (3^o - 2^e)\,a + K.
\]
Since $2^e > 0$, the sign of $b - a$ is determined entirely by the
right-hand side.  We define the \textbf{ordering target function}
\[
  \Delta \;=\; (3^o - 2^e)\,a + K,
\]
so that $\operatorname{sgn}(b - a) = \operatorname{sgn}(\Delta)$.  The two
binary states used to track alignment are:
\[
  s_e \;=\; \operatorname{sgn}(\Delta), \qquad
  s_p \;=\; \operatorname{sgn}(3^o - 2^e),
\]
where $s_e = +$ (resp.\ $-$) means $b > a$ (resp.\ $b \leq a$) and
$s_p = +$ (resp.\ $-$) means $3^o > 2^e$ (resp.\ $3^o \leq 2^e$).

Theorem~\ref{thm:ordering} states that $s_e = s_p$ on every OEEE-free path.
Each step of the path is either an \textbf{OE step} ($v = 1$, $b \equiv 3$
or $7\pmod{8}$, $b_1 = (3b+1)/2$) or an \textbf{OEE step} ($v = 2$,
$b \equiv 1\pmod{8}$, $b_1 = (3b+1)/4$), with carry update $K_1 = 3K + 2^e$
in both cases.  We say an OE step is \textbf{diverging} when $s_p = +$ and
\textbf{converging} otherwise; an OEE step is \textbf{diverging} when
$s_p = -$ and \textbf{converging} otherwise.

\begin{lemma}[Forward sign alignment]
\label{lem:pos}
For any OEEE-free path with $o \geq 1$: if $3^o > 2^e$ then $b > a$.
\end{lemma}
\begin{proof}
Since $3^o$ and $2^e$ are integers, $3^o > 2^e$ implies $3^o - 2^e \geq 1$,
so $\Delta = (3^o - 2^e)a + K \geq a + K \geq a \geq 1 > 0$.
\end{proof}

\subsection{Step dynamics and the $\Delta$ recurrence}

Appending a step with valuation $v \in \{1,2\}$ to a path with parameters
$(o, e, K)$ gives $(o+1, e+v, K_1)$ with $K_1 = 3K + 2^e$.  The updated
target function satisfies:
\begin{align*}
  \Delta_1 &= (3^{o+1} - 2^{e+v})\,a + K_1 \\
           &= 3\bigl[(3^o - 2^e)\,a + K\bigr] + 2^e(3a - 2^v a + 1) \\
           &= 3\Delta + 2^e\,(3 - 2^v)\,a + 2^e.
\end{align*}

\begin{lemma}[Step modifiers]
\label{lem:step}
Under a single step extension:
\begin{enumerate}
  \item[\textup{(OE)}] $\Delta_1 = 3\Delta + 2^e(a+1)$.
  \item[\textup{(OEE)}] $\Delta_1 = 3\Delta - 2^e(a-1)$.
\end{enumerate}
\end{lemma}
\begin{proof}
Substitute $v = 1$ (giving $3 - 2^v = 1$) and $v = 2$ (giving $3 - 2^v = -1$)
into the recurrence above.
\end{proof}

An OE step drives $\Delta$ upward (by $2^e(a+1) > 0$); an OEE step drives it
downward (by $-2^e(a-1) \leq -2^{e+1} < 0$ since $a \geq 3$).

\subsection{Negative parity invariant and main theorem}

A key simplification follows from substituting $K = 2^e b - 3^o a$ (the
path identity) directly into the $\Delta$ recurrence.  For any step:
\[
  K_1 = 3K + 2^e = 3(2^e b - 3^o a) + 2^e = 2^e(3b+1) - 3^{o+1}a,
\]
so
\begin{align*}
  \Delta_1 &= (3^{o+1} - 2^{e+v})a + K_1 \\
           &= (3^{o+1} - 2^{e+v})a + 2^e(3b+1) - 3^{o+1}a \\
           &= 2^e(3b+1) - 2^{e+v}a \\
           &= 2^e\bigl(3b + 1 - 2^v a\bigr).
\end{align*}
Since $2^e > 0$, the sign of $\Delta_1$ is determined entirely by
$3b + 1 - 2^v a$:
\begin{equation}
  \Delta_1 \;\leq\; 0 \;\iff\; 3b + 1 \;\leq\; 2^v a.
  \tag{$\dagger$}\label{eq:delta1}
\end{equation}
For an OEE step ($v = 2$): $\Delta_1 \leq 0 \iff 3b+1 \leq 4a$, i.e.\
$b \leq (4a-1)/3$; for an OE step ($v = 1$): $\Delta_1 \leq 0 \iff
3b+1 \leq 2a$, i.e.\ $b \leq (2a-1)/3$.

We first prove the two OEE cases, then identify precisely what remains open.

\begin{lemma}[OEE negative parity]
\label{lem:neg}
Let $(a, b)$ be an OEEE-free pair with $o \geq 1$.  If the step appended
is an OEE step ($v = 2$) and $3^{o+1} \leq 2^{e+2}$ (new state in
$s_p = -$), then $\Delta_1 < 0$, i.e.\ $b_1 < a$.
\end{lemma}
\begin{proof}
By \eqref{eq:delta1}, $\Delta_1 < 0$ iff $3b + 1 < 4a$.

\textit{Subcase~A — prior state in $s_p = -$ ($3^o \leq 2^e$):}
The inductive hypothesis $\Delta \leq 0$ gives $b \leq a$, so
$3b + 1 \leq 3a + 1$.  Since $a \geq 3$ (odd, $o \geq 1$), $3a + 1 \leq
4a - 2 < 4a$.  Hence $3b + 1 < 4a$ strictly.

\textit{Subcase~B — prior state in $s_p = +$ ($3^o > 2^e$):}
We need $3b + 1 < 4a$.  Substituting $a = (2^e b - K)/3^o$ and using
$K \geq 0$:
\[
  4a \;=\; \frac{4(2^e b - K)}{3^o} \;\geq\; \frac{4 \cdot 2^e b - 4K}{3^o}.
\]
The condition $3b + 1 < 4a$ is equivalent (after rearranging with $a =
(2^e b - K)/3^o$) to $(3b+1) \cdot 3^o < 4(2^e b - K)$, i.e.\
$b(4 \cdot 2^e - 3^{o+1}) > 4K + 3^o$.  Since $3^{o+1} \leq 2^{e+2}$
we have $4 \cdot 2^e - 3^{o+1} = 2^{e+2} - 3^{o+1} \geq 0$.
We prove
\begin{equation}
  b\,(2^{e+2} - 3^{o+1}) \;>\; 4K + 3^o \label{eq:star}\tag{$*$}
\end{equation}
by induction on the number of OEE steps.  \emph{Base} ($o = 0$, $K = 0$):
\eqref{eq:star} gives $b > 1$, which holds for all odd $b \geq 3$.
\emph{Step}: after an OEE step $(e,o,K,b) \mapsto (e+2, o+1, 3K+2^e,
(3b+1)/4)$, using $2^{e+4} - 3^{o+2} = 4(2^{e+2}-3^{o+1}) + 3^{o+1}$:
\begin{align*}
  (3b+1)(2^{e+4} - 3^{o+2})
    &= 4(3b+1)(2^{e+2}-3^{o+1}) + (3b+1)\cdot 3^{o+1} \\
    &> 4\bigl(3(4K+3^o) + 2^{e+2} - 3^{o+1}\bigr) + (3b+1)\cdot 3^{o+1} \\
    &= 48K + 2^{e+4} + (3b+1)\cdot 3^{o+1}.
\end{align*}
The target is $16(3K+2^e) + 4\cdot 3^{o+1} = 48K + 2^{e+4} + 4\cdot
3^{o+1}$, so it suffices that $3b+1 > 4$, i.e.\ $b \geq 3$.  Hence
\eqref{eq:star} holds and $\Delta_1 < 0$.

Both subcases give $\Delta_1 < 0$ whenever the negative parity state is
entered or maintained via an OEE step.
\end{proof}

\begin{remark}
  Subcase~A formalises directly in Lean~4 as a two-line structural induction
  on \texttt{OEF}.  Subcase~B is proved above by an auxiliary induction
  \eqref{eq:star} on the number of OEE steps; this induction is on a
  \emph{different} measure from the structural \texttt{OEF} induction, and no
  clean Lean~4 proof has yet been found.  All numerical evidence confirms the
  subcase is true, and the paper proof is mathematically complete, but the
  Lean formalisation of Lemma~\ref{lem:neg} subcase~B remains open.
\end{remark}

\begin{remark}[Steiner circuit structure and the crossing step]
\label{rem:circuit}
A \textbf{Steiner circuit} is a maximal run from one $b \equiv 1\pmod{8}$
node to the next, consisting of $k \geq 0$ OE steps followed by exactly one
OEE step.  The negative parity state can only be \emph{entered} — as opposed
to sustained — via the terminal OEE step of a circuit.  We call this terminal
step a \textbf{crossing step} when it changes $s_p$ from $+$ to $-$, i.e.\
when $3^o > 2^e$ (before) and $3^{o+1} \leq 2^{e+2}$ (after).

At a crossing step the exponents satisfy the tight sandwich
\[
  2^e \;<\; 3^o \;\leq\; \tfrac{4}{3} \cdot 2^e,
\]
equivalently $3^o/2^e \in (1, 4/3]$.  This forces $4 \cdot 2^e - 3^{o+1}$
to be a small non-negative integer (it lies in $[0, 2^e)$), and the
condition~\eqref{eq:star} becomes a tight numerical inequality between $b$,
$K$, $o$, and $e$ at that single step.

Numerical search ($a \leq 5{,}000$) finds zero violations of $3b+1 < 4a$
among all crossing steps.  The tightest observed margin is
$4a - (3b+1) = 16$ at $(a, b, o, e, K) = (107, 137, 4, 6, 101)$, and the
ratio $3^o/2^e$ at crossing steps lies in $(1.068, 1.266]$, always strictly
below the critical value $4/3$.

The proposed proof strategy for subcase~B is therefore to establish
condition~\eqref{eq:star} via an induction restricted to the OE-steps
\emph{within a single circuit}, starting from the $1\pmod{8}$ circuit-entry
node, rather than a global induction over all OEE steps.  Since each circuit
contains exactly one OEE step, this circuit-level induction collapses
the multi-step auxiliary argument to a single local verification at the
crossing step.
\end{remark}

We can now state precisely what is known and what remains open.

\begin{conjecture}[OE negative parity]
\label{conj:oe}
Let $(a, b)$ be a genuine OEEE-free pair in the negative parity state
($3^o \leq 2^e$) with $b \not\equiv 5\pmod{8}$ and $b \equiv 3$ or
$7\pmod{8}$ (so the next step is an OE step) and $3^{o+1} \leq 2^{e+1}$
(the new state is also $s_p = -$).  Then $3b + 1 \leq 2a$,
equivalently $b \leq (2a-1)/3$.
\end{conjecture}

\noindent\textbf{Why this is not yet proved.}
By \eqref{eq:delta1} with $v = 1$, the condition $\Delta_1 \leq 0$
reduces \emph{exactly} to $3b + 1 \leq 2a$.  We know $b \leq a$ (from
the negative parity state), but $b \leq a$ does not imply $3b+1 \leq 2a$:
the latter requires $b \leq (2a-1)/3$, which is strictly stronger than
$b < a$ for $a \geq 5$ (e.g.\ $a = 7$, $b = 5$ satisfies $b < a$ but
$3 \cdot 5 + 1 = 16 > 14 = 2 \cdot 7$).  The condition $3b+1 \leq 2a$
is a genuine constraint on the orbit — it asserts that when the path is
in the negative parity state and takes an OE step that stays in $s_p=-$,
the node $b$ is bounded by $(2a-1)/3$, not merely by $a$.  This cannot
be deduced from the carry bounds $0 \leq K \leq 4^o - 3^o$ alone:
a phantom predecessor with $b = a - 2$ and $K = 2^e(a-2) - 3^o a$
satisfies the carry bounds and $b < a$, yet violates $3b+1 \leq 2a$
whenever $a \geq 5$.  What is needed is a property of the actual orbit
structure of OEEE-free paths — one that rules out nodes $b$ with
$(2a-1)/3 < b < a$ in the negative parity state.

\begin{theorem}[OEEE-free path ordering, conditional]
\label{thm:ordering}
Let $(a, b)$ be an OEEE-free pair with $o \geq 1$.  Then:
\[
  2^e < 3^o \;\implies\; a < b.
\]
Conditional on Conjecture~\ref{conj:oe}, the full biconditional holds:
\[
  2^e \geq 3^o \;\iff\; a \geq b.
\]
\end{theorem}

\begin{proof}
The forward implication $2^e < 3^o \implies a < b$ is
Theorem~\ref{thm:approx}, unconditional.

For the reverse direction $2^e \geq 3^o \implies a \geq b$, assuming
Conjecture~\ref{conj:oe}: the negative parity state $3^o \leq 2^e$ can
only be entered via an OEE step, for which Lemma~\ref{lem:neg} gives
$\Delta < 0$ strictly.  Conjecture~\ref{conj:oe} then handles all
subsequent OE steps sustaining the negative parity state.  Hence
$\Delta \leq 0$ throughout, giving $b \leq a$.
\end{proof}

The OEEE-free condition is essential: it restricts steps to $v \in \{1,2\}$.
For an OEEE$^+$ step ($v \geq 3$) the modifier becomes
$2^e(3 - 2^v)a + 2^e \leq 2^e(1-5a) < 0$ even when $s_p = +$,
allowing $K$ to escape the parity constraint — as the $a=165$, $b=167$
example illustrates.

See Appendix~\ref{app:empirical} for worked examples illustrating
Theorems~\ref{thm:approx} and~\ref{thm:ordering} and a non-example showing
the necessity of the OEEE-free condition.

\section{Cycles and the $5\pmod{8}$ Constraint}

\begin{theorem}[All-OEE cycle]
\label{thm:nocycle}
The only OEEE-free cycle in which every step is an OEE step ($e = 2o$) is $a = b = 1$.
\end{theorem}

\begin{proof}
Suppose $a = b$ with $e = 2o$.  The path identity gives $K = (4^o - 3^o)a$.
Since $K \leq 4^o - 3^o$, we get $a \leq 1$, hence $a = 1$.  The known
$1\to 1$ cycle ($o=1$, $e=2$) confirms $a = b = 1$ is valid.
\end{proof}

\begin{conjecture}[Unique OEEE-free cycle]
\label{conj:nocycle}
The only cycle on an OEEE-free path is $a = b = 1$.
\end{conjecture}

The all-OEE case is proved by Theorem~\ref{thm:nocycle}.  The mixed case
($o < e < 2o$) requires ruling out cycles with $3^o < 2^e < 4^o$, which
does not follow from the carry bounds alone and sits at
Collatz-conjecture-level difficulty.

\begin{corollary}
\label{cor:oeee-path}
Assuming Conjecture~\ref{conj:nocycle}, if a non-trivial Collatz cycle
exists (other than $1\to 1$), then it must be an OEEE-path.
\end{corollary}

\begin{proof}
By Conjecture~\ref{conj:nocycle}, any OEEE-free cycle is $a = b = 1$
(the all-OEE case is established unconditionally by Theorem~\ref{thm:nocycle};
the mixed case is the conjecture).  Hence any other cycle is not OEEE-free,
i.e.\ it is an OEEE-path.
\end{proof}

\begin{conjecture}
\label{conj:nocycle5}
No non-trivial Collatz cycle contains a $5\pmod{8}$ node.
\end{conjecture}

This conjecture is unproven.  It is offered here not as an established result
but as a suggested avenue for future research: a proof of this conjecture
together with a proof of Conjecture~\ref{conj:nocycle} would establish that
there are no non-trivial Collatz cycles (Theorem~\ref{thm:conditional}).

\begin{theorem}[No non-trivial cycles, conditional]
\label{thm:conditional}
Assuming Conjecture~\ref{conj:nocycle} and Conjecture~\ref{conj:nocycle5},
there are no Collatz cycles other than $1\to 1$.
\end{theorem}

\begin{proof}
By Corollary~\ref{cor:oeee-path} (which assumes Conjecture~\ref{conj:nocycle}),
any non-trivial cycle must be an OEEE-path and hence contains an interior
$5\pmod{8}$ node.  But Conjecture~\ref{conj:nocycle5} asserts no such cycle
exists.  The conclusion follows.
\end{proof}

\section{Conclusion}

The odd Collatz dynamics admit two complementary representations:

\begin{itemize}
  \item A \textbf{proved regular language structure} over $\{1,3,5,7\}$:
        \[
          (7^*3)?(1 \mid 5)
          \quad\text{and}\quad
          \bigl((7^*3)?(1\mid 5)\bigr)^*,
        \]
        encoding Steiner circuits and full orbits.  Here $7$ encodes interior
        OE steps, $3$ encodes the final OE step of a run (immediately before
        the $\mathrm{OEE}^+$ termination), $1$ encodes an OEE
        termination (residue $1 \bmod 8$, exactly two even divisions), and
        $5$ encodes an $\mathrm{OEEE}^+$ termination (residue $5 \bmod 8$,
        three or more even divisions); together $1 \mid 5$ encodes any
        $\mathrm{OEE}^+$ termination (two or more even divisions).  This
        encoding is a theorem following directly from mod-8 arithmetic of $3n+1$.

  \item A \textbf{$5\pmod{8}$ overlay tree} whose nodes are the initial node
        $a$, the final node $b$, all $5\pmod{8}$ nodes on the forward path,
        and for each such $5\pmod{8}$ node $n$, the $0\pmod{3}$ node $m$
        whose forward OEEE-free path reaches $n$, if it exists.
        All edges point in the forward Collatz direction: each $5\pmod{8}$
        node points to the next $5\pmod{8}$ node or to $b$; each $0\pmod{3}$
        node $m$ points to the first $5\pmod{8}$ node on the forward path
        from $m$; $a$ points to the first $5\pmod{8}$ node or directly to $b$.
        By definition, no OEEE-free path visits a $5\pmod{8}$ node.
        Whether every $5\pmod{8}$ node has such a $0\pmod{3}$ node
        remains open (Postulate~\ref{post:halt}).

  \item A \textbf{Collatz path ordering} (Theorems~\ref{thm:approx}
        and~\ref{thm:ordering}): for any two odd integers $a$, $b$ connected
        by a Collatz path with $o \geq 1$, we prove $2^e < 3^o \implies a < b$;
        for OEEE-free pairs the biconditional $2^e \geq 3^o \iff a \geq b$
        is established conditional on Conjecture~\ref{conj:oe}
        (Theorem~\ref{thm:ordering}), via the ordering target function
        $\Delta = (3^o-2^e)a+K$; the OEE cases are proved, and the open
        case is OE steps sustaining the negative parity state.
        All 289 failures of the ordering implication found in adversarial sampling
        occurred on paths with at least one interior $5\pmod{8}$ node, none on
        OEEE-free paths (Table~\ref{tab:empirical}).

  \item A \textbf{cycle constraint}: the only all-OEE cycle is $1\to 1$
        (Theorem~\ref{thm:nocycle}); we conjecture this extends to all
        OEEE-free cycles (Conjecture~\ref{conj:nocycle}), so any other
        Collatz cycle must be an
        OEEE-path.  We conjecture that no Collatz cycle contains a
        $5\pmod{8}$ node (Conjecture~\ref{conj:nocycle5}); conditional on
        this conjecture and Conjecture~\ref{conj:nocycle}, there are no
        Collatz cycles other than $1\to 1$ (Theorem~\ref{thm:conditional}).
\end{itemize}

Together, these results reveal a tight algebraic structure underlying odd
Collatz dynamics.  The residue-$5\pmod{8}$ class plays a distinguished role
throughout: alongside $1\pmod{8}$ it marks the $\mathrm{OEE}^+$ termination
of each Steiner circuit, but uniquely it is the sole source of OEEE$^+$ steps
($v\geq 3$), it determines the branching structure of the overlay tree, and
its absence from OEEE-free path interiors is the key condition under which the
ordering theorem holds.  The empirical evidence (Table~\ref{tab:empirical})
further suggests that interior $5\pmod{8}$ nodes are not merely sufficient
but necessary for ordering failures, pointing to the OEEE-free condition as
the load-bearing structural constraint for the ordering theorem.

\medskip
\noindent\textbf{Open problems.}  One postulate and three conjectures remain
open:

\begin{enumerate}
  \item \textbf{Walk termination} (Postulate~\ref{post:halt}):
        the OEEE-free walk always halts, i.e.\ for every node
        $n \equiv 5\pmod{8}$ the walk eventually reaches a node
        $\equiv 0\pmod{3}$.

  \item \textbf{OE negative parity} (Conjecture~\ref{conj:oe}):
        when an OE step sustains the negative parity state
        ($3^o \leq 2^e$ and $3^{o+1} \leq 2^{e+1}$), the target function
        satisfies $\Delta_1 \leq 0$.  The OEE subcase~A of
        Lemma~\ref{lem:neg} is proved and Lean-formalised; OEE subcase~B
        (entering $s_p=-$ from $s_p=+$) is proved in the paper via
        auxiliary induction~\eqref{eq:star} but the Lean formalisation
        remains open (see the remark after Lemma~\ref{lem:neg}).
        The OE step is the sole remaining gap in Theorem~\ref{thm:ordering}.
        See \S\ref{sec:ordering} for a precise account of why the carry
        bounds alone are insufficient.

  \item \textbf{Unique OEEE-free cycle} (Conjecture~\ref{conj:nocycle}):
        the only cycle on an OEEE-free path is $a = b = 1$.  The all-OEE
        case is proved (Theorem~\ref{thm:nocycle}); the mixed case sits at
        Collatz-conjecture-level difficulty.

  \item \textbf{No $5\pmod{8}$ node in any cycle} (Conjecture~\ref{conj:nocycle5}):
        no Collatz cycle contains a node $\equiv 5\pmod{8}$.
        Conditional on this conjecture and Conjecture~\ref{conj:nocycle},
        there are no Collatz cycles other than $1\to 1$
        (Theorem~\ref{thm:conditional}).
\end{enumerate}

\begin{thebibliography}{9}

\bibitem{erdos1932}
P.~Erd\H{o}s,
\emph{Beweis eines Satzes von Tschebyschef},
Acta Litt. Sci. Szeged \textbf{5} (1932), 194--198.

\end{thebibliography}

\appendix

\section{Empirical Evidence}
\label{app:empirical}

\subsection*{Confirming example: $a=7$, $b=13$}

The forward path $7 \to 22 \to 11 \to 34 \to 17 \to 52 \to 26 \to 13$
has $o = 3$ odd steps (at $7, 11, 17$) and $e = 4$ even halvings
(at $22, 34, 52, 26$), not counting the terminus $13$.
The interior odd nodes are $7, 11, 17$, none $\equiv 5\pmod{8}$, so
$(7, 13)$ is an OEEE-free pair (Definition~\ref{def:branches}).
Then $2^e = 16 < 27 = 3^o$, confirming $2^e < 3^o$
and $a = 7 < 13 = b$, consistent with Theorem~\ref{thm:approx}.

\subsection*{OEEE-path consistent with Theorem~\ref{thm:approx}: $a=55$, $b=71$}

The forward path is
\[
  55 \to 166 \to 83 \to 250 \to 125 \to 376 \to 188 \to 94 \to 47 \to 142 \to 71,
\]
with $o = 4$ odd steps (at $55, 83, 125, 47$) and $e = 6$ even halvings,
not counting the terminus $71$.  The interior node $125 \equiv 5\pmod{8}$,
so $(55, 71)$ is \emph{not} an OEEE-free pair.  Nevertheless,
$2^6 = 64 < 81 = 3^4$, so $2^e < 3^o$, and Theorem~\ref{thm:approx}
applies unconditionally to give $a = 55 < 71 = b$ — as observed.
The OEEE-free hypothesis is irrelevant here: Theorem~\ref{thm:approx}
holds for all Collatz paths.

\subsection*{Confirming example: $a=107$, $b=91$}

The forward path is
\[
  107 \to 322 \to 161 \to 484 \to 242 \to 121 \to 364 \to 182 \to 91,
\]
with $o = 3$ odd steps (at $107, 161, 121$) and $e = 5$ even halvings
(at $322, 484, 242, 364, 182$), not counting the terminus $91$.
The interior odd nodes are $107 \equiv 3\pmod{8}$, $161 \equiv 1\pmod{8}$,
and $121 \equiv 1\pmod{8}$, none $\equiv 5\pmod{8}$, so $(107, 91)$ is an
OEEE-free pair.  Here $2^e = 32 > 27 = 3^o$, so
$2^e \geq 3^o$, and indeed $a = 107 > 91 = b$, consistent with
Theorem~\ref{thm:ordering}.

\subsection*{Non-example: $a=165$, $b=167$}

The forward path is
\[
  165 \to 31 \to 47 \to 71 \to 107 \to 161 \to 121 \to 91 \to 137
      \to 103 \to 155 \to 233 \to 175 \to 263 \to 395 \to 593 \to 445 \to 167,
\]
with $o = 17$ odd steps and $e = 27$ even halvings.

\medskip\noindent\textbf{The OEEE-free condition fails.}
The interior nodes include $165 \equiv 5\pmod{8}$ and $445 \equiv 5\pmod{8}$,
so $(165, 167)$ is \emph{not} an OEEE-free pair and Theorem~\ref{thm:ordering}
does not apply.

\medskip\noindent\textbf{The exponential comparison.}
$2^{27} = 134{,}217{,}728 > 129{,}140{,}163 = 3^{17}$, so $2^e \geq 3^o$
holds — but only just: the excess is $2^e - 3^o = 5{,}077{,}565$, a margin
of less than $4\%$ of $3^o$.

\medskip\noindent\textbf{How close are $b/a$ and $3^o/2^e$?}
From the path identity:
\[
  \frac{b}{a} \;=\; \frac{3^o}{2^e} + \frac{K}{2^e \cdot a}
  \;=\; \frac{129{,}140{,}163}{134{,}217{,}728} + \frac{1{,}106{,}233{,}681}{134{,}217{,}728 \times 165}.
\]
Numerically, $3^o/2^e \approx 0.9622$ while $b/a = 167/165 \approx 1.0121$,
a difference of $K/(2^e a) \approx 0.0499$.  The correction term is large
enough to push $b/a$ above $1$ even though $3^o/2^e < 1$: the path carries
enough extra weight in $K = 1{,}106{,}233{,}681$ to overcome the deficit
$2^e - 3^o = 5{,}077{,}565$.  This is precisely the mechanism that
Theorem~\ref{thm:ordering} asserts cannot occur on an OEEE-free path.

\medskip\noindent\textbf{The admissibility coupling.}
The path identity forces $b \equiv K \cdot (2^e)^{-1} \pmod{3^o}$.
Here $K \cdot (2^{27})^{-1} \bmod 3^{17} = 167 = b$, confirming the
coupling holds — but since $b = 167 \ll 3^{17} = 129{,}140{,}163$, the
residue class places no useful lower bound on $b$.  The OEEE-free condition
is what would otherwise constrain $K$ and force $a \geq b$.

\subsection*{Adversarial sampling}

To provide independent empirical support for Theorem~\ref{thm:ordering}, 2{,}000 random odd starting
points were drawn uniformly from $[1, 10^{12}]$.  For each trajectory the
16 nodes closest in absolute magnitude to the trajectory median were
selected.  This is the most adversarial choice: near the median, $a$ and
$b$ are closest in size, so $b/a \approx 1$ and the margin by which
$b \leq a$ might fail is smallest.  All ancestor-descendant pairs within
each such window were examined (240{,}000 pairs in total, of which 57{,}381
were OEEE-free).

\medskip
\noindent\textbf{What the table shows.}
Table~\ref{tab:empirical} records, for every pair where the ordering
implication $2^e \geq 3^o \implies a \geq b$ fails, how many interior nodes
$\equiv 5\pmod{8}$ the path contains.  The striking feature is the top row:
among the 57{,}381 OEEE-free pairs, \emph{not a single failure} was observed.
Every one of the 289 failures found across all 240{,}000 pairs occurs on a
path containing at least one interior $5\pmod{8}$ node, i.e.\ an OEEE-path.
Moreover the failure count grows with the number of such nodes, suggesting
that each interior $5\pmod{8}$ node contributes independently to inflating
the offset $K$ beyond the level the ordering can absorb.

\begin{table}[h]
\centering
\begin{tabular}{rrc}
\hline
Interior nodes $\equiv 5\pmod{8}$ & Failures & \% of failures \\
\hline
\textbf{0 \quad (OEEE-free paths)} & \textbf{0} & \textbf{0.0\%} \\
1 & 12 & 4.2\% \\
2 & 5  & 1.7\% \\
3 & 133 & 46.0\% \\
4 & 77 & 26.6\% \\
5 & 38 & 13.1\% \\
6 & 21 & 7.3\% \\
7 & 3  & 1.0\% \\
\hline
\textbf{Total} & \textbf{289} & \textbf{100.0\%} \\
\hline
\end{tabular}
\caption{Failures of $2^e \geq 3^o \implies a \geq b$ classified by the
  number of interior nodes $\equiv 5\pmod{8}$ on the path, across 240{,}000
  adversarially-selected ancestor-descendant pairs drawn from 2{,}000 random
  Collatz trajectories with starting points in $[1, 10^{12}]$.  Pairs are
  concentrated near trajectory medians where $b/a \approx 1$, the region
  most likely to produce violations.  Of 57{,}381 OEEE-free pairs examined,
  zero violations were found.  All 289 failures occur on OEEE-paths.}
\label{tab:empirical}
\end{table}

\noindent\textbf{What the table shows.}
The data confirm Theorem~\ref{thm:ordering}: among 57{,}381 OEEE-free pairs,
zero violations of $2^e \geq 3^o \implies a \geq b$ were found.  The table
also illuminates the necessity of the OEEE-free condition: all 289 failures
involve at least one interior $5\pmod{8}$ node.  Notably, many OEEE-paths
satisfy $a \geq b$ even when $2^e \geq 3^o$; the interior $5\pmod{8}$
condition is necessary but not sufficient for failure.

The proof of Theorem~\ref{thm:ordering} is organised via the ordering target
function $\Delta = (3^o - 2^e)a + K$ (Section~\ref{sec:ordering}), whose sign
equals $\operatorname{sgn}(b - a)$.  The negative parity state can only be
entered via OEE steps (Lemma~\ref{lem:neg}), both subcases of which give
$\Delta < 0$ strictly and are fully proved.  OE steps sustaining the
negative parity state are the subject of Conjecture~\ref{conj:oe}, on which
Theorem~\ref{thm:ordering} is conditional.

\end{document}
